\documentclass[12pt,a4paper]{report}

\usepackage{fontspec}
\usepackage{polyglossia}
\setdefaultlanguage{romanian}
\usepackage[margin=2.25cm]{geometry}
\usepackage{setspace}
\usepackage{enumitem}
\usepackage{longtable}
\usepackage{array}
\usepackage{booktabs}
\usepackage{xcolor}
\usepackage{listings}
\usepackage[
    colorlinks,
    linkcolor={black},
    citecolor={black},
    urlcolor={blue}
]{hyperref}

\onehalfspacing
\setlength{\emergencystretch}{3em}
\setlength{\parskip}{0.15em}
\setlist{nosep,leftmargin=*}

\definecolor{codebackground}{rgb}{0.96,0.96,0.96}
\definecolor{codecomment}{rgb}{0.23,0.45,0.30}
\definecolor{codekeyword}{rgb}{0.10,0.20,0.55}
\definecolor{codestring}{rgb}{0.55,0.15,0.15}

\lstdefinestyle{codestyle}{
    backgroundcolor=\color{codebackground},
    basicstyle=\ttfamily\scriptsize,
    breaklines=true,
    frame=single,
    columns=fullflexible,
    keepspaces=true,
    showstringspaces=false,
    numbers=left,
    numberstyle=\tiny\color{gray},
    keywordstyle=\color{codekeyword}\bfseries,
    commentstyle=\color{codecomment},
    stringstyle=\color{codestring},
    tabsize=4
}
\lstset{style=codestyle}

\newcommand{\codechunk}[2]{%
\subsection*{Liniile #1--#2}
\lstinputlisting[language=Python,firstline=#1,lastline=#2,firstnumber=#1]{../src/config.py}
}

\title{Explicarea fisierului \texttt{config.py}}
\author{}
\date{}

\begin{document}

\maketitle
\tableofcontents
\clearpage

\chapter{Rolul fisierului \texttt{config.py}}

Fisierul \path{src/config.py} este centrul de configurare al proiectului. El
defineste constantele globale folosite in pregatirea datelor, procesarea audio,
antrenarea modelelor MaskNet, evaluare, inferenta si rularea experimentelor.
Prin centralizarea acestor valori, restul codului poate importa o singura sursa
de adevar pentru parametri precum sample rate-ul, dimensiunile STFT, directoarele
de rezultate, arhitectura modelului, optimizatorul, loss-ul si preset-urile de
experiment.

Importanta acestui fisier vine din faptul ca multe module depind direct de el:
\texttt{dataset.py} foloseste configuratia datelor, \texttt{mask\_model.py} si
\texttt{train\_mask\_model.py} folosesc configuratia modelului, modulele DSP
folosesc parametrii STFT, iar demo-ul real-time foloseste sample rate-ul,
dimensiunile ferestrelor si directoarele de checkpoint-uri.

\begin{longtable}{p{0.30\textwidth}p{0.60\textwidth}}
\toprule
\textbf{Sectiune} & \textbf{Rol} \\
\midrule
\endhead
Importuri si cai & Stabileste radacina proiectului, directoarele de date,
checkpoint-uri, rezultate si loguri. \\
Audio/STFT & Defineste sample rate-ul, lungimea cadrelor, hop-ul, FFT-ul si
intervalele de lungime audio. \\
Hardware & Alege automat CPU/GPU si seteaza seed-ul random. \\
Dataset & Controleaza split-urile, etichetele VAD, caching-ul si dimensiunile
maxime ale fisierelor. \\
Model & Defineste arhitectura MaskNet si variantele disponibile. \\
Training & Stabileste batch size, numar de epoci, scheduler, early stopping,
AMP si checkpointing. \\
Loss & Alege functia de pierdere si parametrii asociati. \\
Augmentare & Configureaza transformari audio, SpecAugment si Mixup. \\
Preprocesare & Configureaza normalizarea audio si spectrala. \\
Evaluare si inferenta & Controleaza metricile, exemplele salvate si parametrii
de export audio. \\
Preset-uri & Permite aplicarea rapida a unor configuratii standardizate. \\
\bottomrule
\end{longtable}

\chapter{Importuri, directoare si fisiere de log}

\codechunk{1}{23}

\begin{description}
    \item[Liniile 1--3.] Importa \texttt{deepcopy}, \texttt{Path}, \texttt{torch}
    si \texttt{datetime}. \texttt{deepcopy} este folosit mai jos pentru a returna
    configuratii fara ca apelantul sa modifice accidental dictionarele globale.
    \texttt{Path} asigura cai independente de sistemul de operare.

    \item[Liniile 11--15.] Calculeaza radacina proiectului pornind de la locatia
    fisierului curent. \texttt{ROOT} devine directorul \texttt{project\_root}, iar
    directoarele de date sunt construite relativ la acesta. Aceasta decizie face
    proiectul mai portabil, deoarece nu depinde de cai absolute hardcodate.

    \item[Liniile 18--23.] Definesc directoarele pentru checkpoint-uri, rezultate,
    loguri, TensorBoard, ploturi si exemple audio. Aceste constante sunt importate
    de scripturile de training, evaluare si demo pentru a scrie artefactele in
    aceeasi structura.
\end{description}

\codechunk{26}{38}

\begin{description}
    \item[Liniile 26--32.] Functia \texttt{get\_experiment\_dir} creeaza un director
    pentru un experiment. Daca nu se furnizeaza nume, genereaza unul cu timestamp,
    de forma \texttt{experiment\_YYYYMMDD\_HHMMSS}. Astfel, rularile nu se suprascriu
    intre ele.

    \item[Liniile 35--36.] Definesc fisierele de log pentru training si evaluare,
    tot cu timestamp. Acest model ajuta la pastrarea istoricului rularilor.

    \item[Linia 38.] Creeaza directoarele principale daca nu exista. Efectul practic
    este ca restul codului poate salva fisiere fara sa verifice de fiecare data
    existenta directoarelor.
\end{description}

\chapter{Parametrii audio si STFT}

\codechunk{47}{66}

\begin{description}
    \item[Linia 47.] \texttt{SAMPLE\_RATE = 16000} fixeaza frecventa de esantionare
    la 16 kHz. Aceasta valoare este comuna in proiectele de speech enhancement,
    deoarece acopera banda principala a vorbirii si reduce costul computational.

    \item[Liniile 48--53.] Definesc lungimea cadrului si hop-ul in milisecunde,
    apoi le transforma in esantioane. La 16 kHz, 25 ms inseamna 400 esantioane,
    iar 10 ms inseamna 160 esantioane. Aceste valori sunt folosite in STFT, VAD,
    training si inferenta.

    \item[Liniile 56--57.] \texttt{N\_FFT = 512} stabileste dimensiunea FFT-ului,
    iar \texttt{N\_FREQ\_BINS = 257} rezulta din formula \(N/2 + 1\) pentru
    FFT real. Modelul MaskNet lucreaza direct cu aceasta dimensiune spectrala.

    \item[Liniile 60--63.] Definesc lungimea minima si maxima a segmentelor audio
    folosite in dataset. Prin limitarea la 1--10 secunde, training-ul evita fisiere
    prea scurte sau prea lungi.

    \item[Linia 66.] \texttt{SNR\_LEVELS} contine nivelurile standard de mixare
    folosite in pregatirea si evaluarea datelor: 0, 5, 10, 15 si 20 dB.
\end{description}

\chapter{Hardware, seed si reproducibilitate}

\codechunk{73}{83}

\begin{description}
    \item[Liniile 73--74.] Aleg automat dispozitivul de calcul. Daca exista CUDA,
    proiectul foloseste GPU; altfel, foloseste CPU. \texttt{NUM\_GPUS} retine
    numarul de placi grafice disponibile.

    \item[Liniile 77--80.] Seteaza seed-ul random pentru PyTorch. Aceasta masura
    creste reproductibilitatea experimentelor, deoarece initializarile si unele
    operatii random pornesc din acelasi punct.

    \item[Liniile 82--83.] Configureaza cuDNN. \texttt{benchmark=True} permite
    alegerea algoritmilor eficienti pentru dimensiuni de input stabile, iar
    \texttt{deterministic=False} favorizeaza viteza in locul determinismului strict.
\end{description}

\chapter{Configuratia dataset-ului}

\codechunk{90}{116}

\begin{description}
    \item[Liniile 90--95.] \texttt{DATASET\_CONFIG} alege dataset-ul principal,
    aici \texttt{librispeech}, si proportiile pentru train/validation/test:
    80\%, 10\%, 10\%.

    \item[Liniile 98--102.] Controleaza folosirea etichetelor VAD la calculul
    mastilor tinta. \texttt{use\_vad\_labels=False} inseamna ca implicit nu se
    aplica VAD, dar preset-urile pot activa aceasta optiune. \texttt{vad\_soft\_mask\_floor}
    pastreaza o valoare minima pentru mastile soft, astfel incat zonele incerte
    sa nu fie eliminate agresiv.

    \item[Linia 103.] \texttt{on\_the\_fly\_snr\_range} seteaza intervalul SNR
    pentru mixarea dinamica a vorbirii curate cu zgomotul, intre 0 si 20 dB.

    \item[Liniile 105--107.] Configureaza DataLoader-ul. \texttt{num\_workers=0}
    este o alegere importanta pe Windows, deoarece evita problemele de multiprocessing.
    \texttt{pin\_memory=True} poate accelera transferul catre GPU.

    \item[Liniile 110--112.] Controleaza padding-ul, trunchierea fisierelor lungi
    si lungimea maxima in esantioane.

    \item[Liniile 115--116.] Dezactiveaza caching-ul spectrogramelor si audio-ului.
    Aceasta economiseste spatiu pe disc, cu pretul recalcularii unor date.
\end{description}

\chapter{Configuratia modelului MaskNet}

\codechunk{124}{137}

\begin{description}
    \item[Liniile 124--137.] \texttt{MODEL\_CONFIG} este configuratia activa a
    modelului. Ea stabileste numarul de canale de intrare, numarul de canale de
    baza, folosirea LSTM, attention, rezidual, depthwise separable convolution,
    numarul de canale de iesire si tipul mastii.

    \item[Observatie.] Pentru masca de magnitudine, modelul are \texttt{in\_channels=1}
    si \texttt{output\_channels=1}. Pentru masca complexa, preset-ul corespunzator
    va modifica aceste valori la 2, deoarece sunt folosite canalele real si imaginar.
\end{description}

\codechunk{139}{225}

\begin{description}
    \item[Liniile 139--225.] \texttt{MODEL\_VARIANTS} defineste variantele arhitecturale
    disponibile. Fiecare varianta modifica aceeasi configuratie de baza, dar cu
    parametri diferiti.

    \item[\texttt{tiny}.] Varianta usoara, cu \texttt{base\_channels=16}, fara LSTM
    si fara attention. Este potrivita pentru teste rapide.

    \item[\texttt{tiny\_fast}.] Foloseste convolutii depthwise separabile, utile
    pentru antrenare sau inferenta mai rapida pe secvente lungi.

    \item[\texttt{small}, \texttt{balanced}, \texttt{large}, \texttt{xlarge}.]
    Cresc treptat capacitatea modelului. Modelele mai mari pot obtine calitate
    mai buna, dar consuma mai multa memorie si timp de procesare.

    \item[\texttt{balanced\_res}.] Varianta recomandata pentru etapa 2, deoarece
    combina modelul balanced cu blocuri reziduale.

    \item[\texttt{complex\_balanced\_res}.] Varianta pentru Complex MaskNet. Are
    doua canale de intrare si doua de iesire, activare \texttt{tanh} si masca de
    tip \texttt{complex}. Aceasta permite estimarea unei masti complexe real/imaginar.
\end{description}

\codechunk{227}{257}

\begin{description}
    \item[Linia 227.] \texttt{ACTIVE\_MODEL\_VARIANT} stabileste varianta activa la
    importul fisierului. In configuratia curenta, varianta implicita este
    \texttt{tiny}.

    \item[Liniile 230--254.] \texttt{set\_active\_model\_variant} verifica existenta
    variantei cerute si actualizeaza \texttt{MODEL\_CONFIG}. Functia nu creeaza
    modelul direct, ci doar schimba configuratia din care modelul va fi construit.

    \item[Linia 257.] Aplica varianta implicita imediat la import. Astfel,
    \texttt{MODEL\_CONFIG} este mereu sincronizat cu \texttt{ACTIVE\_MODEL\_VARIANT}.
\end{description}

\chapter{Optimizator si preset-uri de optimizare}

\codechunk{261}{323}

\begin{description}
    \item[Liniile 261--278.] \texttt{OPTIMIZER\_CONFIG} contine parametrii generali
    pentru mai multi optimizatori: Adam, SGD si RMSprop. Chiar daca unii parametri
    nu sunt folositi simultan, sunt pastrati centralizat pentru a permite schimbarea
    rapida a optimizatorului.

    \item[Liniile 281--323.] \texttt{OPTIMIZER\_PRESETS} defineste configuratii
    rapide: \texttt{default}, \texttt{aggressive}, \texttt{conservative},
    \texttt{adamw}, \texttt{sgd} si \texttt{rmsprop}.

    \item[Interpretare.] Preset-urile permit comparatii controlate intre strategii
    de optimizare. De exemplu, \texttt{adamw} foloseste weight decay decuplat, iar
    \texttt{conservative} reduce learning rate-ul pentru stabilitate.
\end{description}

\chapter{Configuratia de training}

\codechunk{333}{398}

\begin{description}
    \item[Liniile 333--342.] Seteaza batch size-ul, numarul de epoci, numarul de
    workeri si augmentarea. Valorile sunt orientate catre rulare rapida si
    compatibilitate cu Windows.

    \item[Liniile 344--349.] Definesc learning rate-ul, weight decay-ul si warmup-ul.
    Warmup-ul porneste antrenarea de la un learning rate mic, reducand riscul de
    instabilitate la inceput.

    \item[Liniile 352--369.] Configureaza scheduler-ele disponibile: ReduceLROnPlateau,
    StepLR, CosineAnnealing si OneCycleLR. Nu toate sunt active simultan, dar
    valorile sunt disponibile pentru preset-uri.

    \item[Liniile 372--376.] Configureaza early stopping-ul. Daca validarea nu se
    imbunatateste suficient, antrenarea poate fi oprita.

    \item[Liniile 379--382.] Seteaza gradient clipping-ul, util pentru evitarea
    gradientilor explozivi.

    \item[Liniile 385--386.] Activeaza mixed precision daca exista CUDA. Aceasta
    poate accelera training-ul si reduce consumul de memorie GPU.

    \item[Liniile 389--398.] Controleaza salvarea checkpoint-urilor, logging-ul,
    TensorBoard-ul, validarea si reluarea antrenarii din checkpoint.
\end{description}

\codechunk{402}{467}

\begin{description}
    \item[Liniile 402--467.] \texttt{TRAINING\_PRESETS} defineste moduri de rulare
    pentru diverse scopuri: test rapid, dezvoltare, productie, debug, OneCycleLR
    si configuratia recomandata pentru etapa 2.

    \item[Utilitate.] In loc sa modifici manual multe campuri, un script de training
    poate aplica un preset, iar configuratia devine coerenta imediat.
\end{description}

\chapter{Preset-uri de experiment}

\codechunk{470}{521}

\begin{description}
    \item[Liniile 470--521.] \texttt{EXPERIMENT\_PRESETS} combina preset-uri de
    model, optimizator, training, loss si dataset intr-o singura alegere. Aceasta
    este o forma de configurare la nivel inalt.

    \item[\texttt{stage2\_recommended}.] Alege \texttt{balanced\_res}, \texttt{adamw},
    training \texttt{stage2\_recommended} si loss MSE.

    \item[\texttt{librispeech\_soft\_vad\_recommended}.] Activeaza etichetele VAD
    soft pentru LibriSpeech si foloseste \texttt{balanced\_res}.

    \item[\texttt{librispeech\_complex\_masknet\_recommended}.] Alege varianta
    complexa a modelului, cu VAD soft si AdamW.

    \item[\texttt{librispeech\_fast\_benchmark}.] Alege \texttt{tiny\_fast} si
    reduce lungimea maxima la o secunda, pentru benchmark rapid.
\end{description}

\chapter{Configuratia functiei de pierdere}

\codechunk{524}{586}

\begin{description}
    \item[Liniile 524--538.] \texttt{LOSS\_CONFIG} defineste loss-ul activ si
    parametrii auxiliari. Implicit, loss-ul este \texttt{mse}, dar exista campuri
    pentru loss combinat, Huber, SI-SDR si loss spectral.

    \item[Liniile 542--586.] \texttt{LOSS\_PRESETS} permite alegerea rapida a mai
    multor functii de pierdere: MSE, MAE, Huber, combinat, weighted combined si
    variante mai agresive.

    \item[Interpretare.] MSE penalizeaza mai puternic erorile mari, MAE este mai
    robust la outlieri, Huber combina cele doua comportamente, iar variantele
    weighted pot pune accent suplimentar pe zonele dominante de vorbire.
\end{description}

\chapter{Augmentare si preprocesare}

\codechunk{589}{657}

\begin{description}
    \item[Liniile 589--640.] \texttt{AUGMENTATION\_CONFIG} defineste transformari
    audio: random gain, zgomot, time stretch, pitch shift, reverb, EQ, compresie
    dinamica si clipping. In configuratia de baza, augmentarea globala este
    dezactivata pentru training rapid, dar parametrii sunt pregatiti.

    \item[Liniile 643--650.] \texttt{SPECAUGMENT\_CONFIG} defineste mascari in
    frecventa si timp asupra spectrogramelor. Este dezactivat implicit pentru
    viteza.

    \item[Liniile 653--657.] \texttt{MIXUP\_CONFIG} activeaza Mixup, cu \(\alpha=0.2\)
    si probabilitate 0.3. Mixup combina exemple din acelasi batch pentru
    regularizare.
\end{description}

\codechunk{660}{698}

\begin{description}
    \item[Liniile 660--677.] Configureaza preprocesarea audio: normalizare, eliminare
    DC, preemphasis, AGC si taierea zonelor de liniste.

    \item[Liniile 680--686.] Configureaza normalizarea spectrogramelor. Metoda
    \texttt{per\_channel} cu \texttt{mean\_std} stabilizeaza valorile de intrare
    pentru model.

    \item[Liniile 689--694.] Configureaza conversia in scala logaritmica/dB, cu
    \texttt{top\_db=80}. Aceasta este potrivita pentru spectrograme audio, unde
    dinamica este mare.

    \item[Liniile 697--698.] Activeaza calculul si salvarea statisticilor de
    normalizare.
\end{description}

\chapter{Evaluare, inferenta si experimente}

\codechunk{703}{727}

\begin{description}
    \item[Liniile 703--727.] \texttt{EVAL\_CONFIG} controleaza evaluarea: calculul
    STOI/PESQ, salvarea exemplelor, metricile de training, evaluarea stratificata,
    testele statistice, exemplele audio si spectrogramele.

    \item[Interpretare.] Campul \texttt{metrics\_weights} poate fi folosit pentru
    un scor compozit de selectie a modelelor, unde PESQ, STOI, SNR si scorul
    compozit au ponderi diferite.
\end{description}

\codechunk{736}{797}

\begin{description}
    \item[Liniile 736--762.] \texttt{INFERENCE\_CONFIG} configureaza rularea modelului
    pe fisiere noi: checkpoint, input, output, batch size, dispozitiv, mixed
    precision, normalizare, formatul de iesire si vizualizare.

    \item[Liniile 770--797.] \texttt{EXPERIMENT\_CONFIG} descrie metadatele si
    tracking-ul unui experiment: nume, descriere, tag-uri, versiune de cod,
    informatii sistem, hiperparametri, TensorBoard, W\&B, MLflow si salvarea
    arhitecturii.
\end{description}

\codechunk{805}{823}

\begin{description}
    \item[Liniile 805--823.] \texttt{DISTRIBUTED\_CONFIG} pregateste configuratia
    pentru training distribuit. In proiect este dezactivat implicit, dar campurile
    pentru backend, world size, rank, DataParallel si DistributedDataParallel sunt
    deja definite.
\end{description}

\chapter{Functii de acces si aplicare a preset-urilor}

\codechunk{828}{852}

\begin{description}
    \item[Liniile 828--852.] \texttt{get\_active\_config} intoarce o copie completa
    a configuratiei active. Se foloseste \texttt{deepcopy} pentru dictionare, astfel
    incat apelantul sa nu modifice accidental configuratia globala.

    \item[Utilitate.] Aceasta functie este potrivita pentru salvarea configuratiei
    intr-un checkpoint sau intr-un raport de experiment.
\end{description}

\codechunk{855}{923}

\begin{description}
    \item[Liniile 855--864.] \texttt{apply\_preset} primeste numele preset-ului si
    tipul de configuratie: training, model, optimizer, loss sau experiment.

    \item[Liniile 865--890.] Pentru training, model, optimizer si loss, functia
    verifica daca preset-ul exista, actualizeaza dictionarul corespunzator si
    afiseaza o descriere.

    \item[Liniile 898--919.] Pentru \texttt{experiment}, functia aplica mai multe
    preset-uri in lant: model, optimizer, training, loss si dataset. Acesta este
    mecanismul prin care o singura comanda poate activa o configuratie completa.

    \item[Rol in proiect.] Functia permite rularea experimentelor standardizate
    fara modificarea manuala a fisierului \texttt{config.py}.
\end{description}

\chapter{Afisarea rezumatului configuratiei}

\codechunk{926}{988}

\begin{description}
    \item[Liniile 926--934.] \texttt{print\_config\_summary} afiseaza un header si
    caile principale ale proiectului.

    \item[Liniile 936--940.] Afiseaza dispozitivul, numarul de GPU-uri si starea
    mixed precision.

    \item[Liniile 942--957.] Afiseaza varianta activa a modelului, descrierea,
    canalele de intrare/iesire, canalele de baza, LSTM, attention, rezidual,
    depthwise, tipul mastii si activarea de iesire.

    \item[Liniile 959--966.] Afiseaza parametrii principali de training: batch size,
    epoci, optimizator, learning rate, scheduler si early stopping.

    \item[Liniile 968--977.] Afiseaza informatiile despre dataset si VAD. Daca VAD
    este activ, afiseaza si setul de etichete si podeaua mastii soft.

    \item[Liniile 979--986.] Afiseaza tipul loss-ului si optiunile de evaluare
    principale, apoi inchide raportul cu o linie separatoare.
\end{description}

\chapter{Concluzie}

\texttt{config.py} este un fisier de infrastructura, dar influenteaza aproape
toate componentele proiectului. El stabileste cum sunt citite datele, cum se
construieste modelul, cum se antreneaza, cum se evalueaza si cum ruleaza inferenta.
Prin preset-uri, fisierul permite trecerea rapida intre experimente mici, rulari
de dezvoltare, variante recomandate si modele complexe.

Din perspectiva lucrarii de licenta, acest fisier arata ca proiectul nu este
format doar dintr-un model izolat, ci dintr-un pipeline configurabil: aceleasi
constante STFT sunt folosite in DSP, training si real-time, aceleasi directoare
sunt folosite pentru checkpoint-uri si rezultate, iar aceleasi preset-uri pot
reproduce experimentele prezentate in evaluare.

\end{document}
